\documentclass{article}

\usepackage{graphicx}
\graphicspath{ {./images/} }

\title{Ch. 15 HW \#3 -- Solutions}
\author{Principles of Microeconomics}
\date{\today}

\begin{document}

\maketitle

\section*{Quick Quiz}

\begin{enumerate}

\setcounter{enumi}{5}
\item (d) -- Firms will produce the quantity at which $P = MC$ because they're profit maximizers. They'll enter or exit the market until profits are zero which happens when $P = ATC$.

\item (c) -- Firms will produce the quantity at which $P = MC$ because they're profit maximizers. If new firms want to enter, profits must be positive which means $P > ATC$.

\item (c) -- In the short run, the tax is an increase in fixed costs. Since MC and AVC don't shift, it doesn't change firm's short run quantity produced, so the number of pretzels consumed stays the same. In the long run, firms exit the market because ATC shifts up which makes profit negative. When firms exit, the quantity produced decreases, so the number of pretzels consumed decreases.

\end{enumerate}

\end{document}